% ============================================================
% Converted from post-training-udm-grpo-2026-04-24.md
% ============================================================

\section{UDM-GRPO：均匀离散扩散模型的稳定 GRPO}

\begin{conceptbox}
arXiv: \href{https://arxiv.org/abs/2604.18518}{2604.18518}\\
标题：\textit{UDM-GRPO: Stable and Efficient Group Relative Policy Optimization for Uniform Discrete Diffusion Models}\\
作者：Jiaqi Wang, Haoge Deng, Ting Pan, Yang Liu, Chengyuan Wang, Fan Zhang, Yonggang Qi, Xinlong Wang\\
状态：arXiv v2\\
arXiv 时间：提交于 2026-04-20，更新于 2026-04-21\\
整理日期：2026-04-24\\
相关链接：\href{https://arxiv.org/abs/2604.18518}{arXiv 页面} · \href{https://arxiv.org/html/2604.18518}{arXiv HTML} · \href{https://yovecent.github.io/UDM-GRPO.github.io/}{项目主页} · \href{https://github.com/Yovecent/UDM-GRPO}{代码}
\end{conceptbox}

\subsection{一句话总结}

UDM-GRPO 的核心思想可以记成一句话：

\textbf{不要把“每一步 noisy 的中间预测”当 action，也不要沿模型自己生成的 backward 轨迹做优化；要把“最终干净样本”统一当 action，并且用 forward diffusion 重建训练轨迹。}

这样做以后，GRPO 在 Uniform Discrete Diffusion Model 上终于从“能跑但容易崩”变成了“稳定、有效、还能更高效”。

\subsection{这篇论文到底在讲什么？}

\subsubsection{通俗版}

先把背景讲人话。

这篇论文研究的是：\textbf{怎么把 GRPO 这种强化学习后训练方法，稳定地用到离散扩散图像模型上。}

这里的模型不是连续 latent diffusion，而是\textbf{Uniform Discrete Diffusion Model, UDM}。

它生成图像时，不是在连续空间里一点点修噪声，而是在\textbf{离散 token 空间}里，从“均匀离散噪声”一步步变成最终图像 token。

论文发现，假如你直接把 Flow-GRPO / DDPO 那套思路搬过来，会出现一个很典型的问题：

\begin{itemize}[leftmargin=1.5em]
  \item 前 500 step 左右奖励会上升
  \item 但后面会明显震荡
  \item KL 会爆掉
  \item 最后训练崩掉
\end{itemize}

作者说，这不是“参数没调好”，而是\textbf{建模方式本身就有两个错位}。

\subsubsection{错位 1：把中间预测当 action，不准}

在离散扩散里，模型在早期 timestep 的预测很不稳定、熵很高、很像半成品。

如果你把这些半成品 \code{x\_1\textasciicircum{}t} 当成 RL 的 action 去优化，等于是在强化一个\textbf{本来就不靠谱的中间猜测}。

直觉上就像：

\begin{itemize}[leftmargin=1.5em]
  \item 最终成图是清楚的
  \item 但你非要拿草稿的每一笔都当成“正确动作”去打分
  \item 结果模型学到很多噪声信号
\end{itemize}

所以作者说，\textbf{真正应该被 reward 对齐的，不是每一步的中间猜测，而是最后那张干净图 \code{\textbackslash{}hat\{x\}\_1}}。

\subsubsection{错位 2：沿 backward 轨迹训练，会分布漂移}

预训练时，UDM 学的是：

\begin{itemize}[leftmargin=1.5em]
  \item 从 forward diffusion 生成的 noisy state \code{x\_t}
  \item 去预测最终 clean sample \code{x\_1}
\end{itemize}

但 RL 微调时，如果你直接沿模型自己 rollout 出来的 backward trajectory 优化，那么模型看到的状态分布，已经不是预训练熟悉的那一套了。

这就像：

\begin{itemize}[leftmargin=1.5em]
  \item 预训练时是在标准教材题上练
  \item 微调时突然开始在自己瞎写出来的草稿上继续学
\end{itemize}

于是会出现明显的\textbf{distribution shift / OOD training}，训练自然不稳。

\subsubsection{论文的解决办法}

作者给出两步修正：

\begin{enumerate}[leftmargin=1.5em]
  \item \textbf{动作统一改成最终干净样本}
\end{enumerate}

不管当前在哪个 timestep，都把 action 定义成 \code{\textbackslash{}hat\{x\}\_1}

\begin{enumerate}[leftmargin=1.5em]
  \item \textbf{训练轨迹改成 forward 轨迹}
\end{enumerate}

先正常采样得到最终图像 \code{\textbackslash{}hat\{x\}\_1}，再用 forward diffusion 把它重新扰动成不同 timestep 的 \code{\textbackslash{}hat\{x\}\_t}，然后在这些 \code{\textbackslash{}hat\{x\}\_t} 上做 GRPO

一句话理解：

\textbf{先拿最终样本算 reward，再把这个最终样本“反向做成”一条更接近预训练分布的训练轨迹。}

这就是 UDM-GRPO 稳定的根本原因。

\subsection{先把符号对齐}

这篇论文的时间记号和很多扩散论文不太一样，容易看晕。

\begin{itemize}[leftmargin=1.5em]
  \item \code{x\_1}：最终干净样本
  \item \code{x\_t}：t 时刻的 noisy token 状态
  \item \code{x\_1\textasciicircum{}t}：模型在时刻 \code{t} 对最终样本的中间预测
  \item \code{\textbackslash{}hat\{x\}\_1}：采样得到的最终 clean sample
  \item \code{\textbackslash{}hat\{x\}\_t}：由 \code{\textbackslash{}hat\{x\}\_1} 通过 forward diffusion 扰动得到的训练状态
\end{itemize}

可以粗略记成：

\begin{itemize}[leftmargin=1.5em]
  \item \code{x\_1\textasciicircum{}t} 是“每一步的临时猜测”
  \item \code{\textbackslash{}hat\{x\}\_1} 是“最后真正生成出来的结果”
\end{itemize}

论文的核心就是：

\begin{itemize}[leftmargin=1.5em]
  \item 别对齐 \code{x\_1\textasciicircum{}t}
  \item 去对齐 \code{\textbackslash{}hat\{x\}\_1}
\end{itemize}

\subsection{数学化讲解}

\subsection{1. UDM 的基本形式}

离散扩散工作在离散状态空间

\[
\mathcal{S} = [K]^D
\]

其中：

\begin{itemize}[leftmargin=1.5em]
  \item \code{K} 是 token vocabulary 大小
  \item \code{D} 是 token 序列长度
\end{itemize}

UDM 的起点不是高斯噪声，而是\textbf{均匀离散噪声}：

\[
p_0(x) \sim \mathrm{Unif}([K])^D
\]

目标分布是数据分布 \code{p\_1(x)}。

中间概率路径定义为：

\[
p_t(x) \triangleq \sum_{x_1 \in \mathcal{S}} p_t(x \mid x_1)\, p_1(x_1)
\]

这表示：

\begin{itemize}[leftmargin=1.5em]
  \item 先有一个最终干净样本 \code{x\_1}
  \item 再通过 forward corruption 得到 \code{x\_t}
\end{itemize}

模型训练时做的是“从 noisy state 预测 clean sample”：

\[
\mathcal{L}_{\mathrm{CE}}(\theta)
=
\mathbb{E}_{t, x_1, x_t}\big[-\log p_\theta(x_1 \mid x_t)\big]
\]

这一点很重要，因为后面 UDM-GRPO 的所有改动，本质上都在尽量\textbf{保持和这个预训练目标一致}。

\subsection{2. naive GRPO 是怎么接进来的}

最直接的想法是把 reverse denoising 过程看成一个 MDP。

状态：

\[
s_t \triangleq (c, t, x_t)
\]

在最初探索里，作者把 action 定义成该步对最终图像的预测：

\[
a_t \triangleq x_1^t
\]

对应策略是：

\[
\pi_\theta(a_t \mid s_t)
\triangleq
p_\theta(x_1^t \mid x_t, c)
\]

由于 UDM 输出的是每个 token 的 logits，所以这个概率可以按 token 位置分解：

\[
\pi_\theta(a_t \mid s_t)
=
\prod_{\ell=1}^{D}
\mathrm{Softmax}\!\big(p_\theta(\cdot \mid x_t)_{:,\ell,:}\big)[x_{1,\ell}^{t}]
\]

然后再配上 GRPO 的 group-normalized advantage：

\[
\hat{A}_i
=
\frac{R(x_1^i, c) - \mathrm{mean}(\{R(x_1^i, c)\}_{i=1}^{G})}
{\mathrm{std}(\{R(x_1^i, c)\}_{i=1}^{G})}
\]

再做 clipped policy optimization。

这个思路“形式上可行”，而且确实比原始 URSA 有提升，但很不稳。

\subsection{3. 为什么 naive GRPO 会崩}

论文把原因拆成两个。

\subsubsection{3.1 中间 action 不准确}

早期 timestep 的模型输出熵很高，\code{x\_1\textasciicircum{}t} 还很糊、很 noisy。

但 reward 又只看最终结果。

这就导致一个矛盾：

\begin{itemize}[leftmargin=1.5em]
  \item 奖励来自最终干净样本 \code{\textbackslash{}hat\{x\}\_1}
  \item 更新时却在最大化中间草稿 \code{x\_1\textasciicircum{}t} 的概率
\end{itemize}

于是模型会被迫去模仿很多不可靠的中间猜测，优化方向就被污染了。

\subsubsection{3.2 backward trajectory 和预训练分布不一致}

预训练见到的是：

\[
x_t \sim p_t(x \mid x_1)
\]

也就是\textbf{forward process}产生的 noisy state。

但 RL 微调如果沿 backward rollout 优化，状态来自模型自己生成的 trajectory。

这些状态会逐步偏离预训练分布，形成 OOD。

论文用 FID 和可视化都说明：

\begin{itemize}[leftmargin=1.5em]
  \item \code{\textbackslash{}mathcal\{X\}\_\{forward\}} 比 \code{\textbackslash{}mathcal\{X\}\_\{backward\}} 更接近 \code{\textbackslash{}mathcal\{X\}\_\{pretrain\}}
  \item backward 轨迹在早期步骤偏差尤其大
\end{itemize}

所以训练不稳定并不奇怪。

\subsection{4. UDM-GRPO 的核心改法}

\subsection{4.1 把 action 统一成最终 clean sample}

作者的第一个关键改法是：

\[
a_t \triangleq \hat{x}_1,
\qquad
\pi(a_t \mid s_t)
\triangleq
p_\theta(\hat{x}_1 \mid x_t, c)
\]

直觉上，这样做有两个好处：

\begin{enumerate}[leftmargin=1.5em]
  \item reward 对齐更直接
\end{enumerate}

因为 reward 本来就是打在最终样本上的

\begin{enumerate}[leftmargin=1.5em]
  \item 目标和预训练更一致
\end{enumerate}

因为预训练本来就是“给定 \code{x\_t} 预测 clean sample \code{x\_1}”

也就是说，\textbf{RL 更新不再逼模型去强化噪声很大的半成品，而是强化最终成品本身。}

\subsection{4.2 用 forward process 重建训练轨迹}

第二个关键改法是：

\begin{itemize}[leftmargin=1.5em]
  \item 不沿 backward rollout 的中间状态做优化
  \item 而是先采出 \code{\textbackslash{}hat\{x\}\_1}
  \item 再从 \code{\textbackslash{}hat\{x\}\_1} 出发，用 forward diffusion 重新采样出 \code{\textbackslash{}hat\{x\}\_t}
\end{itemize}

于是训练轨迹变成：

\[
\mathcal{X}_{forward} = \{\hat{x}_t\}_{t=0}^{1},
\qquad
\hat{x}_t \sim p_t(x \mid \hat{x}_1)
\]

这样构造出来的状态分布更接近预训练阶段模型熟悉的输入分布，所以更稳。

\subsection{4.3 重新写成新的 MDP}

把上面两件事合在一起，论文给出新的 MDP：

\[
s_t \triangleq (\hat{x}_t, t, c), \quad
a_t \triangleq \hat{x}_1, \quad
\pi(a_t \mid s_t) \triangleq p_\theta(\hat{x}_1 \mid \hat{x}_t, c),
\]
\[
R(s,a) \triangleq r(\hat{x}_1, c), \quad
\rho_0(s) \triangleq \mathrm{Unif}([K])^D
\]

这其实非常漂亮，因为它把 RL 目标重新压回到了 UDM 的原生预训练语义里：

\begin{itemize}[leftmargin=1.5em]
  \item 输入仍然是 noisy token
  \item 目标仍然是 clean sample
  \item 只是现在 clean sample 的好坏由 reward 来区分，而不是只由监督信号决定
\end{itemize}

\subsection{4.4 policy ratio 怎么变}

对应地，GRPO 里的 policy ratio 也改成：

\[
r_t^i(\theta)
=
\frac{p_\theta(\hat{x}_1^i \mid \hat{x}_t^i, c)}
{p_{\theta_{\mathrm{old}}}(\hat{x}_1^i \mid \hat{x}_t^i, c)}
\]

再把它放进标准的 clipped policy objective：

\[
\mathcal{J}_{policy}^{(t,i)}
=
\min\Big(
r_t^i(\theta)\hat{A}_i,\,
\mathrm{clip}(r_t^i(\theta), 1-\epsilon, 1+\epsilon)\hat{A}_i
\Big)
\]

总目标里再加 reference policy 的 KL 正则。

所以从形式上说，UDM-GRPO 不是发明了一个完全新的 RL loss，

而是\textbf{把 action 和 trajectory 重新定义对了}。

\subsection{5. 两个提速技巧}

\subsection{5.1 Reduced-Step}

如果所有 timestep 都均匀优化，梯度会分散在整条轨迹上，收敛慢。

论文观察到：

\begin{itemize}[leftmargin=1.5em]
  \item 早期 timestep 噪声大、熵高
  \item 也是最容易出错、最值得优化的地方
\end{itemize}

所以他们采用 Reduced-Step：

\begin{itemize}[leftmargin=1.5em]
  \item 对每个样本
  \item 从 diffusion 前半段随机选\textbf{3 个连续 timestep}
  \item 只在这些 timestep 上做 policy optimization
\end{itemize}

这相当于把算力集中到“最关键、最 noisy、最能涨分”的部分。

\subsection{5.2 CFG-Free}

很多 diffusion-RL 方法训练时还带 classifier-free guidance，需要同时处理 conditional 和 unconditional 分支，成本更高。

UDM-GRPO 直接把 CFG 从训练里拿掉。

直觉上：

\begin{itemize}[leftmargin=1.5em]
  \item 一开始样本质量可能会更差
  \item 但训练过程更简单
  \item 也减少了不必要的计算和优化耦合
\end{itemize}

论文结果显示，CFG-Free 后期反而更强，尤其 OCR 上更明显。

\subsection{训练流程速记}

把论文算法压缩成一句可执行的话：

\begin{enumerate}[leftmargin=1.5em]
  \item 给定 prompt，先用旧策略采样 \code{G} 个最终样本 \code{\textbackslash{}hat\{x\}\_1}
  \item 用 reward 给这 \code{G} 个样本打分，算出 group-normalized advantage
  \item 对每个最终样本，选几个早期 timestep
  \item 用 forward diffusion 从 \code{\textbackslash{}hat\{x\}\_1} 采出对应的 \code{\textbackslash{}hat\{x\}\_t}
  \item 计算 \code{p\_\textbackslash{}theta(\textbackslash{}hat\{x\}\_1\ \textbackslash{}mid\ \textbackslash{}hat\{x\}\_t,\ c)} 和 ratio
  \item 用 GRPO clipped loss + KL 去更新模型
\end{enumerate}

所以它本质上是：

\textbf{sample final image -> score final image -> rebuild forward states -> train on those states}

\subsection{实验设置速记}

\begin{itemize}[leftmargin=1.5em]
  \item 基座模型：URSA 的 1.7B text-to-image model
  \item 训练任务：组合生成、OCR 文本渲染、人类偏好对齐
  \item 数据与 reward：沿用 Flow-GRPO 的三套设置
  \item group sampling：每个 batch 是 \code{16\ groups\ x\ 8\ samples}
  \item 优化器：AdamW
  \item 学习率：\code{1e-6}
  \item 训练采样步数：默认 \code{10} 步
  \item 评估步数：默认 \code{25} 步
  \item 训练硬件：\code{32\ x\ A100\ 40GB}
\end{itemize}

\subsection{主要实验结果}

\subsection{1. headline numbers}

论文摘要和项目页给出的 headline 是：

\begin{itemize}[leftmargin=1.5em]
  \item GenEval：\code{69\%\ ->\ 96\%}
  \item PickScore：\code{20.46\ ->\ 23.81}
  \item OCR：\code{8\%\ ->\ 57\%}
\end{itemize}

这里要注意一个小细节：

这三个 headline 数字来自不同对照行，\textbf{并不是严格同一个 baseline 行里的三列同时对比}。所以读摘要时别把它理解成“完全同一基线三项同时从 A 到 B”。

\subsection{2. 主表里真正能记的结果}

Table 2 里最值得记的是：

\begin{itemize}[leftmargin=1.5em]
  \item \code{UDM-GRPO}: GenEval \code{0.96}, PickScore \code{23.81}, OCR \code{0.57}
  \item \code{URSA}: GenEval \code{0.69}, PickScore \code{21.79}, OCR \code{0.08}
  \item \code{URSA\ (w/o\ CFG)}: GenEval \code{0.36}, PickScore \code{20.46}, OCR \code{0.04}
\end{itemize}

和强连续模型对比：

\begin{itemize}[leftmargin=1.5em]
  \item \code{SD3.5-L}: GenEval \code{0.71}, PickScore \code{22.91}, OCR \code{0.68}
  \item \code{FLUX.1-Dev}: GenEval \code{0.66}, PickScore \code{22.84}, OCR \code{0.59}
\end{itemize}

\subsubsection{这说明什么？}

\begin{enumerate}[leftmargin=1.5em]
  \item \textbf{在 GenEval 上非常强}
\end{enumerate}

\code{0.96} 对 \code{0.69} 是一个巨大跃升，而且明显高于 SD3.5-L / FLUX / URSA 本体

\begin{enumerate}[leftmargin=1.5em]
  \item \textbf{在 PickScore 上也很强}
\end{enumerate}

\code{23.81} 超过表中的强连续基线 \code{22.91\ /\ 22.84}

\begin{enumerate}[leftmargin=1.5em]
  \item \textbf{OCR 提升巨大，但不是全表最强}
\end{enumerate}

相对 URSA 的 \code{0.08} 确实是大幅提升到 \code{0.57}，但表里 \code{SD3.5-L\ =\ 0.68}，\code{FLUX\ =\ 0.59}，所以 OCR 上更准确的说法是：

\textbf{对离散 UDM 基座提升极大，但还没绝对超过最强连续模型。}

这个细节我觉得很重要，因为它说明：

\begin{itemize}[leftmargin=1.5em]
  \item UDM-GRPO 确实解决了“离散扩散没法稳定做 RL”的问题
  \item 但在“文字渲染”这件事上，问题还没完全结束
\end{itemize}

\subsection{消融最值得记的结论}

Table 3 非常关键，因为它直接告诉你每个设计到底值不值。

\subsubsection{1. 最原始接法：backward + \code{x\_1\textasciicircum{}t}}

\begin{itemize}[leftmargin=1.5em]
  \item GenEval \code{0.84}
  \item PickScore \code{21.99}
  \item OCR \code{0.23}
\end{itemize}

说明：

直接把中间预测当 action 虽然比原始 URSA 强，但还远不够好。

\subsubsection{2. 只改 action：backward + \code{\textbackslash{}hat\{x\}\_1}}

\begin{itemize}[leftmargin=1.5em]
  \item GenEval \code{0.89}
  \item PickScore \code{23.10}
  \item OCR \code{0.23}
\end{itemize}

说明：

只把 action 改成最终 clean sample，就已经明显提升了，尤其是 GenEval 和 PickScore。

这支持了论文的第一个洞见：\textbf{final clean sample 比 intermediate prediction 更适合做 action。}

\subsubsection{3. 再改 trajectory：forward + \code{\textbackslash{}hat\{x\}\_1}}

\begin{itemize}[leftmargin=1.5em]
  \item GenEval \code{0.94}
  \item PickScore \code{23.51}
  \item OCR \code{0.34}
\end{itemize}

说明：

forward 轨迹确实更稳、更强。

这支持第二个洞见：\textbf{轨迹分布要贴近预训练分布。}

\subsubsection{4. 最终完整版本：forward + \code{\textbackslash{}hat\{x\}\_1} + CFG-Free}

\begin{itemize}[leftmargin=1.5em]
  \item GenEval \code{0.96}
  \item PickScore \code{23.81}
  \item OCR \code{0.57}
\end{itemize}

说明：

CFG-Free 不是“省点算力的小技巧”，而是最终性能的一部分。

\subsection{Reduced-Step 的结论}

论文还比较了三种 timestep 优化策略：

\begin{itemize}[leftmargin=1.5em]
  \item 只优化早期高噪声 timestep
  \item 随机连续 timestep
  \item 所有 timestep 都优化
\end{itemize}

结论是：

\begin{itemize}[leftmargin=1.5em]
  \item 在 GenEval 和 PickScore 上，\textbf{优化早期高噪声 timestep 最好}
  \item 在 OCR 上差距没那么大
\end{itemize}

这和直觉是一致的：

\begin{itemize}[leftmargin=1.5em]
  \item 组合关系、布局、属性绑定，很多是在早期决定的
  \item 后期更多像修纹理和补细节
\end{itemize}

\subsection{我对这篇论文的理解}

我觉得这篇最有价值的地方，不是又写了一个“diffusion + RL”的损失函数，而是它把一个很容易混乱的问题讲清楚了：

\textbf{在离散扩散里，RL 到底应该对齐什么对象，在哪条轨迹上对齐。}

作者的答案非常干净：

\begin{itemize}[leftmargin=1.5em]
  \item 对齐对象：最终 clean sample
  \item 对齐轨迹：forward diffusion trajectory
\end{itemize}

这两个选择都不是拍脑袋，而是分别对应两个原则：

\begin{enumerate}[leftmargin=1.5em]
  \item \textbf{reward consistency}
\end{enumerate}

reward 打在最终样本上，就该让最终样本做 action

\begin{enumerate}[leftmargin=1.5em]
  \item \textbf{pretraining consistency}
\end{enumerate}

预训练是在 forward-noised state 上学 clean prediction，RL 也应该尽量待在这个分布附近

所以我会把这篇论文记成：

\textbf{它不是在发明更花哨的 GRPO，而是在把 UDM 的 RL 问题重新表述对。}

\subsection{一种最好记的记忆口诀}

如果只记一句话：

\textbf{UDM-GRPO = final sample as action + forward process as trajectory + early-step optimization + CFG-Free training}

如果再展开一点：

\begin{itemize}[leftmargin=1.5em]
  \item \code{action} 选最终图，不选中间草稿
  \item \code{trajectory} 选 forward，不选 backward rollout
  \item \code{timestep} 多优化早期高噪声部分
  \item \code{training} 去掉 CFG，后期更稳更快
\end{itemize}

\subsection{这篇论文的局限}

\subsubsection{1. OCR 还没有绝对统治}

它把 URSA 的 OCR 从 \code{0.08} 拉到 \code{0.57} 已经很夸张了，但相比 \code{SD3.5-L\ =\ 0.68} 仍有差距。

所以“稳定 RL for UDM”解决了，但“离散扩散的文字生成天花板”还没有彻底解决。

\subsubsection{2. 主要创新是重构 RL 表述，不是基础模型本身}

也就是说，这篇的价值更偏向：

\begin{itemize}[leftmargin=1.5em]
  \item post-training formulation
  \item RL stability
  \item trajectory design
\end{itemize}

而不是“发明了一个更强的新 UDM backbone”。

\subsubsection{3. 对更复杂 reward / 更长采样链的泛化还可以继续看}

论文附录里展示了对 FUDOKI 也有效，这很好；

但如果以后 reward 更复杂、采样更长、模型更大，这套“forward rebuild + final action”的优势还能保持多少，还值得继续观察。

\subsection{我会怎么把它和别的工作区分开}

和很多 diffusion RL 论文相比，UDM-GRPO 最不一样的点在于：

\begin{itemize}[leftmargin=1.5em]
  \item 别人经常在想“loss 怎么改”
  \item 它更像在想“MDP 应该怎么定义才不别扭”
\end{itemize}

这是一种更底层、也更值得借鉴的思路。

如果以后你看到别的“离散生成 + RL”工作，其实都可以拿这篇的框架去问两个问题：

\begin{enumerate}[leftmargin=1.5em]
  \item 你优化的 action，真的是 reward 最关心的对象吗？
  \item 你训练时看到的 state distribution，和预训练时一致吗？
\end{enumerate}

很多训练不稳定的问题，本质上都能被这两个问题捅出来。

\subsection{参考链接}

\begin{itemize}[leftmargin=1.5em]
  \item arXiv Abstract: \url{https://arxiv.org/abs/2604.18518}
  \item arXiv HTML: \url{https://arxiv.org/html/2604.18518}
  \item Project Page: \url{https://yovecent.github.io/UDM-GRPO.github.io/}
  \item Code: \url{https://github.com/Yovecent/UDM-GRPO}
\end{itemize}
